\section{Visión del mundo: determinismo, libre albedrío y el sentido de la vida}

\subsection{Un proceso de Márkov determinista}

En la parte filosófica de la entrevista, Weng Jiayi (翁家翌) mostró una faceta muy distinta de la de los temas técnicos. Sostuvo con firmeza:

\begin{importantbox}{La visión del mundo de Weng Jiayi: determinismo}
\textbf{Vivimos dentro de un proceso de Márkov determinista.} Todo puede predecirse: lo que piensas, la siguiente palabra que vas a decir, todo quedó fijado en el instante mismo del Big Bang. El ser humano no tiene libre albedrío, y el mundo macroscópico no tira los dados (la aleatoriedad de la mecánica cuántica a nivel microscópico puede explicarse como una «modificación de las líneas de mundo en segundo plano»).
\end{importantbox}

Él reconoce que se trata de una «visión del mundo bastante pesimista» y que, en el fondo, tampoco quisiera aceptarla, como si él mismo se hubiera convertido en un átomo simulado. Pero, con base en su experiencia personal, considera que en efecto es un hecho. Afirma que «ya lo he verificado innumerables veces», aunque no puede revelar los detalles de esa experiencia personal.

El entrevistador insistió: ¿y la incertidumbre de la mecánica cuántica? La respuesta de Weng Jiayi fue que «a nivel macroscópico no se tiran los dados, a nivel microscópico sí»; pero la aleatoriedad microscópica puede explicarse como una «modificación de algunas líneas de mundo en segundo plano». Reconoció que «se puede pensar que lo que digo son puras tonterías», pero aun así mantuvo su postura.

Esta visión del mundo genera una tensión interesante con su trabajo en RL: el supuesto central de RL es que el agente puede modificar la recompensa futura eligiendo distintas acciones, pero si el mundo es determinista, entonces la «elección» misma del agente también estaría predeterminada. ¿Es Weng Jiayi consciente de esta tensión? Su respuesta lo sugiere: «hay quien quiere desarrollar este tipo de modelos (que predicen el futuro) solo para averiguar cuál es la regularidad que rige este mundo por debajo».

\subsection{La paradoja de invertir en el futuro}

El entrevistador señaló con agudeza una contradicción en la visión del mundo de Weng Jiayi: si todo está determinado, se llegará a ese punto tanto si inviertes en el futuro como si no, entonces ¿para qué invertir?

La respuesta de Weng Jiayi fue: «de todos modos es mejor invertir un poco». Pero al insistírsele en que «bajo un proceso determinista tu inversión no tendría ningún efecto», reconoció: «invertir en el futuro también podría ser algo determinista; el que inviertas o no tampoco es tu libre albedrío».

\subsection{Un periodo de desorientación vital}

Al final de la entrevista, Weng Jiayi compartió con franqueza su estado actual:

\begin{knowledgebox}{El estado actual de Weng Jiayi}
Weng Jiayi atraviesa un periodo de desorientación en cierto punto de su carrera. El trabajo en RL Infra, que alguna vez le gustó mucho, se ha ido volviendo cada vez más determinista con el paso del tiempo. Dice que «alguna vez tuve claro qué quería, pero ahora ya no lo tengo claro». Su meta a corto plazo es retirarse anticipadamente, acumular capital suficiente y luego dedicar tiempo a encontrar lo que de verdad quiere hacer. Respecto a sí mismo dentro de diez años, su única expectativa es «hacer lo que quiera hacer en ese momento, con los recursos y la capacidad suficientes».
\end{knowledgebox}

\subsection{Resumen de la sección}

Las reflexiones filosóficas de Weng Jiayi le dan profundidad a toda la entrevista. Un ingeniero que construyó la infraestructura central de OpenAI resulta, en el terreno filosófico, un determinista absoluto. Considera que la AGI es «un hecho consumado» y que todo lo demás es «algo bastante determinista», que «ya se le ve el final». Pero, paradójicamente, esa certeza es justo lo que lo hace sentirse desorientado: si todo está determinado, ¿dónde queda el sentido del esfuerzo?

Respecto a sí mismo dentro de diez años tiene una única expectativa: «hacer lo que quiera hacer en ese momento, con los recursos y la capacidad suficientes». El entrevistador insistió: «¿no vas a intervenir en lo que él quiera en ese momento?». Weng Jiayi respondió: «porque las ideas cambian; puede que no importe lo que tú quieras». Y añadió: «lo único que puedo hacer ahora es invertir en ese yo de entonces, o dicho de otro modo, invertir en el futuro, para que él tenga la libertad de elegir».

Su mensaje final fue: \textbf{«explorar, al final de cuentas, qué es lo que uno mismo quiere: esta pregunta vale la pena pensarla toda la vida».}

